% ===================================================================
%  اللوحة نمرة ٣ — الطفولة والشباب.
%
%  Traduction, faite en 2026, en arabe egyptien ecrit, sur l'ido de
%  texto/io. Regles de la colonne : voir 10-lawha-01.tex. Deux
%  scenes, deux numerotations : les cles portent c1 et c2.
%
%  LES JOURS DE LA SEMAINE SONT EGYPTIENS, ET ILS NE RESSEMBLENT PAS
%  A CEUX DE L'ARABE STANDARD :
%
%    (يوم) الاتنين   lundi     MSA : الإثنين
%    (يوم) التلات    mardi     MSA : الثلاثاء
%    (يوم) الأربع    mercredi  MSA : الأربعاء
%
%  C'est la meme regle phonologique que celle du tableau 2 — le ث
%  passe a ت — appliquee cette fois au calendrier, et le nom du jour
%  s'en trouve raccourci d'une syllabe.
%
%  LE VETEMENT DE FETE PARLE ITALIEN, ET C'EST L'EGYPTE DU XIXe
%  SIECLE QUI L'A FAIT ENTRER :
%
%    برنيطة (50, 59, 74)  le chapeau   MSA : قبعة
%    جوانتي (51)          les gants    MSA : قفازات
%    كرافتة (9)           la cravate
%    صديري (17)           le gilet
%    بنطلون (16)          le pantalon
%    مايوه (3)            le maillot de bain
%    كاسكيت (11)          la casquette
%    شُرّاب (6)            les chaussettes  MSA : جوارب
%    قبقاب (75)           le sabot
%
%  LA CEREMONIE DU BAPTEME TOMBE JUSTE, MOT POUR MOT. L'egyptien
%  chretien a exactement les deux fonctions du livret : الإشبين (49)
%  le parrain, الإشبينة (53) la marraine ; et ce qu'on jette aux
%  enfants sur le parvis s'appelle الملبّس (55), les dragees, qui est
%  precisement ce que la planche montre.
%
%  ET LA FETE FORAINE EST UN مولد (39). Ce n'est pas une traduction
%  approchee : le مولد egyptien est la fete patronale d'un saint, avec
%  ses manèges, son cirque, sa ménagerie et ses baraques, et c'est le
%  seul mot qui nomme la chose entiere. De meme, le maire (37) du
%  bourg est العمدة, qui est le titre exact du chef d'un village
%  egyptien.
%
%  DEUX AUTRES MOTS DE CE TABLEAU MERITENT LEUR LIGNE : بتاعة اللبن
%  (24) la laitiere — l'egyptien nomme le marchand par بتاع + sa
%  marchandise, et c'est le meme بتاع qui sert a la possession ; et
%  بيعيّط (40) « il pleure », la ou l'arabe standard dit يبكي.
%
%  L'ARABE N'A DEMANDE AUCUNE INVERSION DE RENVOI, et ce n'est pas un
%  hasard : comme l'ido, il place le possede AVANT le possesseur —
%  « سهم البرج », la fleche du clocher. Le chinois faisait l'inverse
%  et coutait trente-cinq reprises a la colonne cantonaise ; ici
%  l'annexion suit l'ido pas a pas.
% ===================================================================

%%K t03-apar-1 apar
\VUsaut{3.0mm}
\VUtitre{15.0pt}{60}{اللوحة نمرة ٣}
\VUsaut{1.6mm}
\VUtitre{11.4pt}{0}{الطفولة والشباب.}
\VUsaut{3.4mm}

%%K t03-apar-2 apar
(اللوحتين رقم ٣ ورقم ٤ متجمّعين كده علشان يقدّموا في أربع \VUgras{مشاهد
عائلية} الأفكار الأساسية عن \VUgras{الجو}، و\VUgras{السن}، و\VUgras{العيلة}،
و\VUgras{الهدوم}، و\VUgras{الحالة}.)

\VUsaut{4.0mm}
%%K t03-tit-1 sub
\VUcentre{12.0pt}{0}{\textit{المشهد الأول.}}
\VUsaut{3.0mm}

%%K t03-tit-2 sub
\VUpk{10.2pt}{40}{حفلة معمودية في قرية.}
\VUsaut{2.4mm}

%%K t03-tit-3 sub
\VUpk{10.2pt}{40}{الربيع.}
\VUsaut{3.0mm}

%%K t03-c1-01-1 p
1. --- إحنا في \VUgras{الربيع}، في شهر \VUgras{مارس}، أو \VUgras{أبريل}، أو
\VUgras{مايو}، في يوم من أيام \VUgras{الأسبوع}: \VUgras{الاتنين}، أو
\VUgras{التلات}، أو \VUgras{الأربع}. والساعة \VUgras{عشرة} الصبح.

%%K t03-c1-02-1 p
2. --- \VUgras{الجو} حلو، و\VUgras{الهوا} صافي. والشمس منوّرة. وفي كام
\VUgras{سحابة صغيّرة} \textsuperscript{(2)} طالعة في \VUgras{السما} الزرقا
\textsuperscript{(1)}.

%%K t03-c1-03-1 p
3. --- \VUgras{الشجر} لسه ما فيهوش \VUgras{ورق} كتير. \VUgras{شجر الحور}
\textsuperscript{(3)} الرفيّع و\VUgras{الدلب} \textsuperscript{(17)} العالي لسه ما
عندهمش غير براعم.

%%K t03-c1-04-1 p
4. --- \VUgras{عصافير السنونو} \textsuperscript{(4)} رجعت وبتلفّ وهي بتزقزق من
الفرحة حوالين \VUgras{السهم} \textsuperscript{(7)} بتاع \VUgras{برج الأجراس}
\textsuperscript{(5)} اللي متوّج \VUgras{الكنيسة} \textsuperscript{(6)} بتاعة القرية.

%%K t03-tit-4 sub
\VUsaut{3.4mm}
\VUpk{10.2pt}{40}{الكنيسة.}
\VUsaut{3.0mm}

%%K t03-c1-05-1 p
5. --- الكنيسة دي بسيطة أوي، بـ\VUgras{بوّابتها} \textsuperscript{(13)} الكبيرة
و\VUgras{ساعتها} \textsuperscript{(8)} التخينة. \VUgras{العقرب الصغير}
\textsuperscript{(9)} على رقم X، وده بيدلّ على \VUgras{الساعات}.

و\VUgras{العقرب الكبير} \textsuperscript{(10)} على XII (نُصّ)، وده بيدلّ على
\VUgras{الدقايق}.

%%K t03-c1-06-1 p
6. --- وفي البرج، \VUgras{الأجراس} \textsuperscript{(11)} بتدقّ من الفرحة. وفوق
السقف واقف \VUgras{صليب} \textsuperscript{(12)} صغير من حجر.

%%K t03-c1-07-1 p
7. --- الكنيسة مش عندها \VUgras{بوّابة} \textsuperscript{(13)} كبيرة بس، عندها كمان
باب صغير جنبي. ومن الباب الصغير ده بيخرج حضرة \VUgras{الخوري}
\textsuperscript{(15)} لابس \VUgras{جُبّته}، بيسيب \VUgras{حجرة الملابس}
\textsuperscript{(14)} ورايح \VUgras{بيت الخوري} \textsuperscript{(19)}.

%%K t03-c1-08-1 p
8. --- وجنبه، أهي \VUgras{بتاعة اللبن} \textsuperscript{(24)} مع \VUgras{حمارها}
\textsuperscript{(23)}. راجعة من المدينة، اللي ودّت لها لبن كويس للعيال.

%%K t03-tit-5 sub
\VUsaut{4.0mm}
\VUpk{10.2pt}{40}{جنينة الفاكهة.}
\VUsaut{3.4mm}

%%K t03-c1-09-1 p
9. --- وسيادة الحمار مبسوط أوي من الجرية الصبحية دي. بيشمّ بمزاجه الهوا اللي
عطّرته ريحة \VUgras{الورد} \textsuperscript{(20)} اللي مغطّي \VUgras{شجر الفاكهة}
\textsuperscript{(18)} في \VUgras{جنينة الفاكهة} اللي جنبهم. \VUgras{شجر الخوخ}
\textsuperscript{(18)} و\VUgras{شجر المشمش} \textsuperscript{(19)} كله وردي وأبيض،
وبيبشّر بمحصول كويس.

%%K t03-c1-10-1 p
10. --- وفي نفس الوقت، \VUgras{النحل} \textsuperscript{(22)} الصغيّر بتاع
\VUgras{الخلية} \textsuperscript{(21)} رايح هناك يلمّ \VUgras{عسله} الحلو وهو بيطنّ.

%%K t03-c1-11-1 p
11. --- بس في إيه؟ أهم عيال القرية جايين بيجروا وبيهرّبوا \VUgras{الوز}
\textsuperscript{(25)} اللي اتخضّ. ده الناس بتخرج من الكنيسة بعد \VUgras{معمودية}
بنت الموثّق.

%%K t03-c1-12-1 p
12. --- ويعقوب الصغير، في \VUgras{سريره} \textsuperscript{(26)}، صحي من دقّ الأجراس
الفرحان، وأخته مجدلينا، وهي كمان عايزة تشوف موكب المعمودية، بتترجّى أمها
تيجي تمشّط لها شعرها وتلبّسها هدومها على عتبة البيت.

%%K t03-c1-13-1 p
13. --- والأم وافقت. لبست \VUgras{طرحتها} \textsuperscript{(28)}،
و\VUgras{صدريتها} \textsuperscript{(29)}، و\VUgras{مريلتها} \textsuperscript{(30)}، وجت
قعدت على كرسي صغير قدام الباب. ومجدلينا لسه ما عليهاش غير
\VUgras{جيبتها التحتية} \textsuperscript{(31)} و\VUgras{كورسيهها} \textsuperscript{(32)}.

%%K t03-c1-14-1 p
14. --- ما أسعدها! نسيت الورد اللي بتحبه، \VUgras{القرنفل} \textsuperscript{(34)}
بتاعها اللي \VUgras{الفراشات} \textsuperscript{(35)} بتحطّ عليه، و\VUgras{التوليب}
\textsuperscript{(36)} بتاعها، و\VUgras{زنبق الوادي} \textsuperscript{(37)} بتاعها، اللي
في \VUgras{القصارى} \textsuperscript{(38)} على \VUgras{الحيطة} \textsuperscript{(33)}،
و\VUgras{الورد البلدي} \textsuperscript{(39)} بتاعها اللي عامل سياج حلو أوي. هي قاعدة
تتفرّج على هدوم ستّات الموكب الغنية.

%%K t03-c1-15-1 p
15. --- عيال القرية مش واخدين بالهم أوي من الهدوم. بيتزاحموا ويخبطوا في بعض
علشان ياخدوا \VUgras{ملبّس} \textsuperscript{(55)} أو \VUgras{قروش}
\textsuperscript{(52)}. وواحد فيهم بيعيّط \textsuperscript{(40)}.

%%K t03-c1-16-1 p
16. --- وواحد تاني داس على رجل \VUgras{الكلب} \textsuperscript{(42)}، اللي جري وهو
بينبح وهرّب \VUgras{الفرخة} \textsuperscript{(43)} و\VUgras{كتاكيتها}
\textsuperscript{(44)}.

%%K t03-tit-6 sub
\VUsaut{5.0mm}
\VUpk{10.2pt}{40}{موكب المعمودية.}
\VUsaut{4.0mm}

%%K t03-c1-17-1 p
17. --- قدام الموكب ماشية \VUgras{الداية} \textsuperscript{(45)}. لابسة
\VUgras{بونيه} \textsuperscript{(46)} حلو بشرايط طويلة. وشايلة \VUgras{المولود}
\textsuperscript{(47)} كله ملبّس \VUgras{دانتيلا} \textsuperscript{(48)}.

%%K t03-c1-18-1 p
18. --- ووراها جايين \VUgras{الإشبين} \textsuperscript{(49)} و\VUgras{الإشبينة}
\textsuperscript{(53)}. وهما بيوزّعوا الملبّس والقروش. الإشبين لابس
\VUgras{جوانتي} \textsuperscript{(51)} و\VUgras{برنيطة عالية} \textsuperscript{(50)}.
والإشبينة ماسكة بإيدها \VUgras{بوكيه} \textsuperscript{(54)} حلو.

%%K t03-c1-19-1 p
19. --- وبعديهم بنشوف \VUgras{عم} \textsuperscript{(56)} الطفل و\VUgras{عمّته}
\textsuperscript{(58)}. العمّة لابسة \VUgras{برنيطة} \textsuperscript{(59)} شيك، والعم
لابس \VUgras{ردنجوت} \textsuperscript{(57)} أسود وصديري أبيض. وأهو بعدهم \VUgras{ابن
العم} \textsuperscript{(60)} و\VUgras{بنت العم} \textsuperscript{(61)}. ودي ماسكة بإيدها
\VUgras{أخت} \textsuperscript{(62)} المولود، برتا الصغيرة.

%%K t03-c1-20-1 p
20. --- و\VUgras{أخو} \textsuperscript{(63)} برتا التاني ماشي ورا. وهو ماسك إيد
\VUgras{قريبة} \textsuperscript{(64)} ليهم. و\VUgras{أصحاب} \textsuperscript{(65)} العيلة
جايين بعد كده.

%%K t03-tit-7 sub
\VUsaut{4.6mm}
\VUpk{10.2pt}{40}{المتفرّجين.}
\VUsaut{3.4mm}

%%K t03-c1-21-1 p
21. --- وجنب الموكب، أهو \VUgras{شحّات} \textsuperscript{(66)} متّكي على
\VUgras{عكازه} \textsuperscript{(67)}. بيطلب صدقة.

%%K t03-c1-22-1 p
22. --- وعلى الناحية التانية في ولد صغير \VUgras{مؤدّب} \textsuperscript{(68)} أوي.
بيسلّم ويقول «\VUgras{صباح الخير}» للناس اللي يعرفهم.

%%K t03-c1-23-1 p
23. --- وأهل القرية خرجوا من بيوتهم علشان يشوفوا موكب المعمودية. بنشوف
\VUgras{فلّاح} \textsuperscript{(69)} لابس \VUgras{طاقيته القطن} وجنبه
\VUgras{فلّاحة} \textsuperscript{(70)}. وبعدين \VUgras{عجوزة} \textsuperscript{(71)} متّكية
على \VUgras{عصايتها} \textsuperscript{(72)}. وأخيرًا \VUgras{الطحّان} \textsuperscript{(73)}
بـ\VUgras{برنيطة اللبد} \textsuperscript{(74)} الكبيرة بتاعته، وفلّاحة صغيرة لابسة
\VUgras{قباقيب} \textsuperscript{(75)}، وبتعلّم ابنها الصغير يمشي.

%%K t03-c1-24-1 p
24. --- المشهد ده مسلّي أوي.

\VUsaut{4.0mm}
%%K t03-tit-8 sub
\VUcentre{12.0pt}{0}{\textit{المشهد التاني.}}
\VUsaut{3.0mm}

%%K t03-tit-9 sub
\VUpk{10.2pt}{40}{عيد عام.}
\VUsaut{2.4mm}

%%K t03-tit-10 sub
\VUpk{10.2pt}{40}{ألعاب المية والسباق.}
\VUsaut{3.0mm}

%%K t03-c2-01-1 p
1. --- \VUgras{الصيف} جه. وشمس شهر \VUgras{يونيو} (يوليو أو \VUgras{أغسطس})
الحامية بتشجّع الشباب إنهم يستحمّوا ويجدّفوا.

%%K t03-c2-02-1 p
2. --- على الشطّ اليمين بتاع النهر، المجدّفين بيقلعوا هدومهم علشان يشتركوا في
ألعاب المية والسباق.

%%K t03-c2-03-1 p
3. --- واحد فيهم بيقلع \VUgras{جزمته} \textsuperscript{(1)}؛ وبيفكّ رباطها. وزميليه
جاهزين خلاص. مش فاضل عليهم غير \VUgras{الفانلة} \textsuperscript{(2)}
و\VUgras{المايوه} \textsuperscript{(3)}. وبيتكلّموا وهما مستنّيين أصحابهم. بيتكلّموا
عن المولد والعيد اللي على الشطّ التاني.

%%K t03-c2-04-1 p
4. --- وعلى الأرض، جنبهم، مرمية هدوم زمايلهم: جوز \VUgras{بوت}
\textsuperscript{(4)}، و\VUgras{جزم} \textsuperscript{(5)}، و\VUgras{شُرّاب}
\textsuperscript{(6)}، وبرانيط \VUgras{لبد} \textsuperscript{(7)} و\VUgras{قش}
\textsuperscript{(8)}، و\VUgras{كرافتة} \textsuperscript{(9)}.

%%K t03-c2-05-1 p
5. --- وواحد من الشباب طوى \VUgras{بدلته} \textsuperscript{(10)} بعناية. وحطّها على
فوطة، اللي زميله هيلفّ فيها البدلتين حالًا.

%%K t03-c2-06-1 p
6. --- والولد السادس اتأخّر. لسه حاطط \VUgras{الكاسكيت} \textsuperscript{(11)} على
راسه. وما قلعش لا \VUgras{قميصه} \textsuperscript{(12)}، ولا \VUgras{ياقته}
\textsuperscript{(13)}، ولا \VUgras{مانشيتاته} \textsuperscript{(14)}. وماسك بإيده
\VUgras{صديريه} \textsuperscript{(17)}، ولسه لابس \VUgras{بنطلونه} \textsuperscript{(16)}
و\VUgras{حمّالاته} \textsuperscript{(15)}. أكيد مش هيبقى جاهز للسباق اللي جاي.

%%K t03-c2-07-1 p
7. --- وجنب الشباب في طريق صغير على جوانبه \VUgras{شوك} \textsuperscript{(18)} كتير،
بيودّي لحافّة النهر، اللي فيها عمّ يعقوب بيحضّر بين \VUgras{البوص}
\textsuperscript{(19)} \VUgras{الألعاب النارية} \textsuperscript{(20)} اللي هيولّعوها
بالليل.

%%K t03-tit-11 sub
\VUsaut{4.4mm}
\VUpk{10.2pt}{40}{السباق.}
\VUsaut{3.4mm}

%%K t03-c2-08-1 p
8. --- وعلى النهر، كام ستّ، وهما مستخبّيين تحت شمسياتهم، بيتمشّوا
بـ\VUgras{مركب} \textsuperscript{(21)}. وبيتابعوا السباق اللي مشوّق أوي. في كل
\VUgras{قارب} \textsuperscript{(22)}، أربع \VUgras{مجدّفين} \textsuperscript{(24)} بيجدّفوا
بكل قوّتهم،

%%K t03-c2-08-1 p suite
وفي نفس الوقت \VUgras{الدفّاف} \textsuperscript{(26)} ماسك \VUgras{الدفّة}
\textsuperscript{(25)} وبيوجّه المركب الرفيّع. و\VUgras{المجاديف} \textsuperscript{(23)}
بتضرب المية على وزن واحد، والمراكب الخفيفة بتجري بسرعة تدوّخ.

%%K t03-c2-09-1 p
9. --- وكام شاب بيتنافسوا، جنب عمود الكوبري، على جايزة \VUgras{عمود السارية}
\textsuperscript{(28)}. واحد فيهم بيمشي بالراحة على العمود اللي بيتلوّى وبيزحلق
علشان يوصل لـ\VUgras{العلم} \textsuperscript{(29)} الصغير المربوط في آخره.
و\VUgras{منافسه} \textsuperscript{(30)} قليل البخت وقع في المية. وبيعوم علشان يطلع بر.

%%K t03-c2-10-1 p
10. --- والشباب الأكبر \VUgras{بيتصارعوا على المراكب} \textsuperscript{(32)} جنب
\VUgras{الرصيف} \textsuperscript{(33)}. والألعاب دي كلها بيتفرّج عليها \VUgras{جمهور}
\textsuperscript{(31)} كبير، متّكي على سور \VUgras{الكوبري} \textsuperscript{(27)}
ومتكوّم على الشطّ. بس في متفرّجين بيلفّوا علشان يتابعوا لعبة \VUgras{عمود
الجايزة} \textsuperscript{(45)}، أو علشان يتفرّجوا على \VUgras{موكب العمدة} اللي جاي
علشان \VUgras{يوزّع} الجوايز.

%%K t03-c2-11-1 p
11. --- الأول، أهم كام \VUgras{عيّل} \textsuperscript{(35)} ماشيين قدام
\VUgras{الموسيقيين} \textsuperscript{(36)}؛ وبعدين جاي \VUgras{العمدة} \textsuperscript{(37)}،
والمساعدين، و\VUgras{المجلس البلدي} بتاع البلد. وأخيرًا ماشيين
\VUgras{المطافي} \textsuperscript{(38)}.

%%K t03-tit-12 sub
\VUsaut{5.0mm}
\VUpk{10.2pt}{40}{المولد.}
\VUsaut{3.6mm}

%%K t03-c2-12-1 p
12. --- في ناس كتير أوي في \VUgras{ساحة المولد} \textsuperscript{(39)}. الناس بتيجي
تتفرّج على \VUgras{الخيام} المختلفة: \VUgras{الخيل الخشب} \textsuperscript{(40)}،
و\VUgras{السيرك} \textsuperscript{(41)} اللي العرض ابتدا فيه خلاص؛ و\VUgras{الرقص
العام} \textsuperscript{(42)}، اللي فيه شباب وبنات البلد بيستمتعوا بمزاج
\VUgras{الرقص}.

%%K t03-c2-13-1 p
13. --- وفي الآخر، بنشوف \VUgras{الحلبة} \textsuperscript{(44)}، اللي فيها
\VUgras{المصارع} الإسباني الشجاع بيصارع

%%K t03-c2-13-1 p suite
\VUgras{التور} \textsuperscript{(45)} الهايج؛ وبعدين \VUgras{الزحليقة}
\textsuperscript{(49)}، و\VUgras{ميدان الرماية} \textsuperscript{(46)} اللي الناس فيه
بتتمرّن على \VUgras{البنادق} وعلى الضرب على \VUgras{الهدف}.

%%K t03-c2-14-1 p
14. --- وجنبه، أهو \VUgras{مسرح المولد} \textsuperscript{(47)} اللي بيمثّلوا فيه
\VUgras{الكوميديا} و\VUgras{الدراما}. وقريّب أوي منه، خيمة \VUgras{العملاق}
\textsuperscript{(48)} و\VUgras{القزم}. \VUgras{طول} الأول مترين وخمستاشر سنتيمتر؛
والتاني بالكاد يبقى قد عيّل.

%%K t03-c2-15-1 p
15. --- وأخيرًا، أهو \VUgras{معرض الوحوش} \textsuperscript{(50)} الكبير، بـ\VUgras{حيواناته
الوحشة}، و\VUgras{النمر} \textsuperscript{(51)} بـ\VUgras{ضوافره} القوية،
و\VUgras{الأسد} \textsuperscript{(52)} بـ\VUgras{لبدته} الكتيفة، و\VUgras{الزرافتين}
\textsuperscript{(53)}، و\VUgras{الفيل} \textsuperscript{(54)} اللي بيستعمل
\VUgras{خرطومه} بمهارة كبيرة.

%%K t03-c2-16-1 p
16. --- و\VUgras{المروّض} \textsuperscript{(55)} اللي بيدرّب الحيوانات دي واقف على
باب الخيمة.

\VUsaut{6.0mm}
\VUfilet{22mm}
